\chapter{Operations \& debugging}
\label{ch:operations-debugging}

This chapter documents where logs live on disk, how to control their
verbosity, and the commands you typically need when a user reports
odd behavior and you want to see what the app actually did.

\section{Where the logs live}
\label{sec:logs-location}

Starting with \textbf{v1.7.0}, mybibli writes its logs to two places
simultaneously:

\begin{itemize}
  \item \textbf{stdout} (unchanged) — visible via
        \texttt{docker compose logs mybibli}. Best for live-tailing
        during an active session.
  \item \textbf{A daily-rotating file on disk}, by default under
        \texttt{/var/log/mybibli/} inside the container, mounted on a
        named volume (\texttt{mybibli\_logs}) on the host so files
        survive \texttt{docker compose pull \&\& up -d}. Best for
        \emph{post-mortem} debugging.
\end{itemize}

Files are named \texttt{mybibli.log.YYYY-MM-DD} (one per UTC day).
A background task purges files older than 30 days automatically.

\subsection*{Changing the in-container log path}
\label{sec:log-dir-override}

The default in-container log directory \texttt{/var/log/mybibli/} is
controlled by the \texttt{MYBIBLI\_LOG\_DIR} environment variable.
Change it only if you remount the volume elsewhere — the shipped
\texttt{docker-compose.yml} maps the \texttt{mybibli\_logs} named
volume (or a bind mount via \texttt{LOGS\_HOST\_PATH}) onto
\texttt{/var/log/mybibli}, so changing one without the other lands
logs in the writable layer where they get wiped on every
\texttt{docker compose up -d} after a \texttt{pull}.

See section~\ref{sec:install-logs-bind} for the bind-mount setup
(host-visible logs for DSM Log Center / journald shipping).

\section{Tailing logs from the host}

\begin{lstlisting}
# Tail the current day's file in real time
docker compose exec mybibli tail -f /var/log/mybibli/mybibli.log.$(date -u +%Y-%m-%d)

# Last 100 lines of the current day's file
docker compose exec mybibli tail -n 100 /var/log/mybibli/mybibli.log.$(date -u +%Y-%m-%d)

# Show all available days
docker compose exec mybibli ls /var/log/mybibli/

# Search across all retained days for a keyword
docker compose exec mybibli sh -c "grep ERROR /var/log/mybibli/mybibli.log.*"
\end{lstlisting}

If you mounted the volume as a bind mount (\texttt{LOGS\_HOST\_PATH}
in \texttt{.env}, see section~\ref{sec:install-logs-bind}), the files
are visible directly on the host without needing
\texttt{docker compose exec}.

\section{Log levels}
\label{sec:log-levels}

Two environment variables control verbosity:

\begin{itemize}
  \item \texttt{MYBIBLI\_LOG\_LEVEL} (preferred, v1.7.0+) — accepts
        either a plain level (\texttt{trace}, \texttt{debug},
        \texttt{info}, \texttt{warn}, \texttt{error}) or a
        \texttt{tracing-subscriber} \texttt{EnvFilter} directive
        list, e.g. \texttt{mybibli=debug,sqlx::query=warn}.
  \item \texttt{RUST\_LOG} — legacy fallback, same format. Honored
        when \texttt{MYBIBLI\_LOG\_LEVEL} is unset.
\end{itemize}

Default is \texttt{info} — production-safe (info + warn + error,
no per-query SQL noise, no per-request trace floods).

\subsection*{Recommended levels}

\begin{description}
  \item[\texttt{info}] Production default. Catches authentication
        events, metadata fetches, scheduled-task runs, user-facing
        errors. Roughly 10--100 lines per day on a single-household
        install.
  \item[\texttt{mybibli=debug}] Active investigation of an mybibli
        feature. Adds request-routing, cache hits/misses, scan
        dispatch decisions. Roughly 10$\times$ the info volume.
  \item[\texttt{mybibli=trace}] Deepest one-shot investigation.
        Includes per-call tracing within mybibli's own code. Very
        noisy; use briefly and bump back to \texttt{info}.
  \item[\texttt{mybibli=debug,sqlx::query=info}] Active investigation
        AND see every SQL statement (useful when a query is suspect).
\end{description}

\subsection*{How to change the level}

\textbf{Two paths}, depending on whether you want the change to
survive container restarts or just for one investigation.

\subsubsection*{From the admin UI (v1.7.1+, recommended)}

Log in as Admin, go to \texttt{/admin?tab=system}, scroll to the
\textbf{Logging} section. Enter the new directive in the
\texttt{Log level} field (anything that
\texttt{tracing-subscriber::EnvFilter} accepts \string— plain
level or full directive list). Click \textbf{Save}. The next log
line uses the new filter \string— no container restart, no
\texttt{docker compose up -d}. The value is persisted to the
\texttt{settings} table, so it survives container restarts.

\subsubsection*{From \texttt{docker-compose.yml} / \texttt{.env}}

Edit the file, set \texttt{MYBIBLI\_LOG\_LEVEL}, then:

\begin{lstlisting}
docker compose up -d
\end{lstlisting}

The container restarts and the new level takes effect from the next
log line. \textbf{Caveat:} the admin-saved value in the
\texttt{settings} table wins on the next boot, so if you've ever
saved a level via the admin UI, that's what will load \string—
the env var only matters on a fresh install before the first save.

\section{Reading the structured JSON}

Logs are emitted as line-delimited JSON. Each line carries:

\begin{itemize}
  \item \texttt{timestamp} — ISO~8601 UTC.
  \item \texttt{level} — \texttt{TRACE} / \texttt{DEBUG} /
        \texttt{INFO} / \texttt{WARN} / \texttt{ERROR}.
  \item \texttt{target} — the Rust module that emitted the line
        (e.g. \texttt{mybibli::routes::catalog}). Filter on this to
        narrow a search to a subsystem.
  \item \texttt{fields.message} — the human-readable summary.
  \item additional \texttt{fields.*} — contextual keys like
        \texttt{user\_id}, \texttt{volume\_id}, \texttt{isbn},
        \texttt{title\_id}.
\end{itemize}

To pretty-print:

\begin{lstlisting}
docker compose exec mybibli tail -n 50 /var/log/mybibli/mybibli.log.$(date -u +%Y-%m-%d) \
  | jq -r '[.timestamp, .level, .fields.message] | @tsv'
\end{lstlisting}

\section{Sharing a log excerpt for a bug report}

When opening an issue at \url{https://github.com/guycorbaz/mybibli/issues}:

\begin{enumerate}
  \item Note the wall-clock time when the bug occurred (e.g.
        ``around 14:30 UTC today'').
  \item Extract the relevant 5-minute window:
  \begin{lstlisting}
  # Replace 14:25 / 14:35 with your window.
  docker compose exec mybibli sh -c \
    "awk '/\"timestamp\":\"$(date -u +%Y-%m-%d)T14:2[5-9]/,/\"timestamp\":\"$(date -u +%Y-%m-%d)T14:3[0-5]/' \
       /var/log/mybibli/mybibli.log.$(date -u +%Y-%m-%d)"
  \end{lstlisting}
  \item Paste the output into the issue body, fenced as
        \texttt{```json}.
\end{enumerate}

Redact anything you consider sensitive — log lines may contain
borrower names, ISBNs of titles you own, or session-token prefixes
(but never full session tokens or password hashes; those are kept
out of log emission by design).

\section{Retention \& rotation}

\begin{itemize}
  \item Rotation happens \textbf{automatically} at UTC midnight,
        handled in-process by \texttt{tracing-appender::rolling}.
  \item Files older than \textbf{30 days} are purged once per day by
        a background task. The retention window is hardcoded in
        v1.7.0; a future release will expose it as an admin setting.
  \item The current day's file is never deleted (the cutoff is
        strict).
\end{itemize}

If you want to keep logs longer than 30 days, copy the rotated files
out of the volume to long-term storage:

\begin{lstlisting}
docker compose cp mybibli:/var/log/mybibli/. ./logs-archive/
\end{lstlisting}
